\begin{table}[htbp]
\centering
\scriptsize
\setlength{\tabcolsep}{2.5pt}
\caption{Size and power of the nine back-tests measured by the validator of case C22: the rejection rate at the 5 percent level when the null holds at its boundary and against one stated alternative, over two thousand repetitions each, with the Wilson 95 percent interval in brackets; the size of a test that rejects at its red or Monte Carlo level is read at that level.}
\label{tab:market-sizepower}
\begin{tabular}{p{2.9cm}p{6.6cm}rr}
\toprule
Test & Sample, level and alternative & Size (\%) & Power (\%) \\
\midrule
\relax McNeil-Frey residual test & 2,500 days at 5\%, 125 exceedances expected; the ES forecast is 95\% of the true 5\% expected shortfall of a normal loss (mean residual 0.103) & 4.5 [3.6, 5.4] & 93.2 [92.0, 94.2] \\
\relax Kupiec proportion of failures & 1,500 days at 5\%; the VaR is exceeded 1.3 times too often (a hit rate of 6.5\%, not 5\%) & 5.1 [4.2, 6.2] & 69.3 [67.2, 71.3] \\
\relax Christoffersen independence & 1,500 days at 5\%; a hit is followed by a hit with probability 15\%, not 5\% (clustered exceedances at the same hit rate) & 6.6 [5.6, 7.8] & 87.5 [86.0, 88.9] \\
\relax Christoffersen conditional coverage & 1,500 days at 5\%; a hit is followed by a hit with probability 15\%, not 5\% (clustered exceedances at the same hit rate) & 6.7 [5.6, 7.8] & 81.2 [79.4, 82.9] \\
\relax Duration test & 400 days at 5\%; a hit is followed by a hit with probability 40\%, not 5\% (clustered exceedances at the same hit rate) & 5.4 [4.5, 6.5] & 63.8 [61.7, 65.9] \\
\relax Acerbi-Szekely $Z_2$ & 250 days at 2.5\%; the realised losses are 1.25 times as dispersed as the VaR and ES forecasts say & 4.8 [3.9, 5.8] & 89.4 [88.0, 90.7] \\
\relax Du-Escanciano, unconditional & 2,000 days at 10\%; the lower tail has too much mass (u = r\textasciicircum{}1.05: 11.2\% of the PITs below 0.10, not 10\%) & 4.7 [3.9, 5.7] & 56.5 [54.3, 58.7] \\
\relax Du-Escanciano, conditional & 2,000 days at 10\%; the tail PITs (below 0.10) cluster: one follows another with probability 20\%, not 10\% & 5.5 [4.6, 6.6] & 73.2 [71.2, 75.1] \\
\relax Dynamic quantile test & 1,500 days at 5\%; a hit is followed by a hit with probability 15\%, not 5\% (clustered exceedances at the same hit rate) & 4.8 [3.9, 5.8] & 73.7 [71.7, 75.6] \\
\relax Dynamic quantile test & 250 days at 1\%; a hit is followed by a hit with probability 10\%, not 1\% (clustered exceedances at the same hit rate) & 5.1 [4.2, 6.1] & 16.0 [14.4, 17.6] \\
\relax Dynamic quantile test & 250 days at 2.5\%; a hit is followed by a hit with probability 15\%, not 2.5\% (clustered exceedances at the same hit rate) & 5.9 [5.0, 7.1] & 25.7 [23.8, 27.7] \\
\relax Dynamic quantile test & 750 days at 1\%; a hit is followed by a hit with probability 10\%, not 1\% (clustered exceedances at the same hit rate) & 4.2 [3.4, 5.1] & 27.9 [26.0, 29.9] \\
\bottomrule
\end{tabular}
\end{table}
